\section{Till They Find Rest in Thee}

Augustine began the \emph{Confessions} with a sentence addressed to God:

\begin{quote}
Thou movest us to delight in praising Thee; for Thou hast formed us for Thyself, and our hearts are restless till they find rest in Thee.\endnote{Augustine, \emph{Confessions} I.1, trans.\ J.~G. Pilkington, \emph{Nicene and Post-Nicene Fathers}, 1st series, vol.~1 (Buffalo: Christian Literature Publishing Company, 1886).}
\end{quote}

Restlessness of this kind belongs to a creature made for an end it does not yet possess, and it is the first sign of heaven that each of us carries. Thomas Aquinas, expounding the last article of the Creed, cites this very sentence when he comes to say what everlasting life is. It is, he says, first of all union with God, and that union consists in perfect vision, in the highest praise, and in the perfect satisfying of desire; for nothing created ever satisfies the desire of man: \latin{Deus enim solus satiat, et in infinitum excedit}, God alone satisfies, and exceeds it without end. That is why, he goes on, the Lord says \latin{intra in gaudium domini tui}; and he adds a saying that he gives as Augustine's, though it is Anselm's, from the \emph{Proslogion}, and that turns the sentence round: \latin{totum gaudium non intrabit in gaudentes, sed toti gaudentes intrabunt in gaudium}. The whole of that joy will not enter into those who rejoice; those who rejoice will enter, whole, into the joy.\endnote{Thomas Aquinas, \emph{Collationes in Symbolum Apostolorum}, a.~12, Latin text of the Corpus Thomisticum. The English renderings in the text are glosses, not a published translation. The saying is Anselm's: \emph{Proslogion} 26 (Latin).}

\subsection{Ostia}

Once, for a moment, Augustine had been given something of this joy. Not many days before his mother fell sick of her last fever, the two of them stood at a window in a house at Ostia, looking into the garden, resting before a voyage. He tells what they spoke of:

\begin{quote}
We then were conversing alone very pleasantly; and, ``forgetting those things which are behind, and reaching forth unto those things which are before,'' we were seeking between ourselves in the presence of the Truth, which Thou art, of what nature the eternal life of the saints would be, which eye hath not seen, nor ear heard, neither hath entered into the heart of man.\endnote{Augustine, \emph{Confessions} IX.10.23, trans.\ Pilkington, NPNF 1.1; the onset of Monica's fever, IX.11.27.}
\end{quote}

Their talk climbed. They passed through all bodily things, ``and even the heaven itself, whence sun, and moon, and stars shine upon the earth,'' and came to their own minds, and went beyond them, toward that Wisdom by whom all things are made; ``and while we were thus speaking, and straining after her, we slightly touched her with the whole effort of our heart; and we sighed, and there left bound `the first-fruits of the Spirit;' and returned to the noise of our own mouth, where the word uttered has both beginning and end.''\endnote{Augustine, \emph{Confessions} IX.10.24, trans.\ Pilkington, NPNF 1.1.} Then, still speaking together, mother and son reasoned toward heaven. ``We were saying, then,'' Augustine writes:

\begin{quote}
If to any man the tumult of the flesh were silenced,---silenced the phantasies of earth, waters, and air,---silenced, too, the poles; yea, the very soul be silenced to herself, and go beyond herself by not thinking of herself,---silenced fancies and imaginary revelations, every tongue, and every sign, and whatsoever exists by passing away, \ldots\ and He alone speak not by them, but by Himself, that we may hear His word, not by fleshly tongue, nor angelic voice, nor sound of thunder, nor the obscurity of a similitude, but might hear Him---Him whom in these we love---without these, \ldots\ If this could be sustained, and other visions of a far different kind be withdrawn, and this one ravish, and absorb, and envelope its beholder amid these inward joys, so that his life might be eternally like that one moment of knowledge which we now sighed after, were not this ``Enter thou into the joy of Thy Lord''?\endnote{Augustine, \emph{Confessions} IX.10.25, trans.\ Pilkington, NPNF 1.1.}
\end{quote}

Augustine's passage works by subtraction, as the Apocalypse does: everything that passes is silenced, and every created voice falls still, until God speaks ``not by them, but by Himself.'' And the words they found for what would then be given are the Lord's own, from the parable of the talents: ``Enter thou into the joy of thy lord.''\endnote{Matt 25:21, 23, Douay--Rheims (Challoner). In the Postconciliar Missal the parable is the Gospel of the Thirty-third Sunday in Ordinary Time, Year A: \emph{Ordo lectionum Missae}, editio typica altera (1981), Matt 25:14--30. Monica's death soon after is told in \emph{The Debt of Nature}.}

Thomas explains why the parable says it that way round. Joy is full when nothing remains to be desired, and in heaven nothing will remain, ``so that the joy of the blessed is full to perfection---indeed over-full, since they will obtain more than they were capable of desiring.'' And since no creature can hold the joy that God deserves, ``it follows that this perfectly full joy is not taken into man, but, on the contrary, man enters into it, according to Matt.\ 25:21: `Enter into the joy of thy Lord.'\,''\endnote{Aquinas, \emph{Summa theologiae} II-II, q.~28, a.~3, trans.\ Fathers of the English Dominican Province (New York: Benziger Brothers; Project Gutenberg eBook 18755).}

\subsection{Rest without weariness}

There is a fear that comes to Christians when they think of an endless happiness: that it will be the same thing for ever, and so at last a burden. Benedict XVI named the fear in the encyclical on hope: the word eternal, he wrote, ``suggests to us the idea of something interminable, and this frightens us.'' His answer was that eternity ``is not an unending succession of days in the calendar,'' but ``a plunging ever anew into the vastness of being, in which we are simply overwhelmed with joy.''\endnote{Benedict XVI, encyclical \emph{Spe salvi} (30 November 2007), n.~12, official English text.} Fathers and saints had answered it long before, each with an image of his own.

Gregory of Nyssa, in his dialogue \emph{On the Soul and the Resurrection}, gives to his sister, whom he calls the Teacher, the teaching that souls were made as vessels for the blessings of God, and that the vessel grows with what it receives:

\begin{quote}
The All-creating Wisdom fashioned these souls, these receptacles with free wills, as vessels as it were, for this very purpose, that there should be some capacities able to receive His blessings and become continually larger with the inpouring of the stream. Such are the wonders that the participation in the Divine blessings works: it makes him into whom they come larger and more capacious; from his capacity to receive it gets for the receiver an actual increase in bulk as well, and he never stops enlarging. The fountain of blessings wells up unceasingly\ldots\endnote{Gregory of Nyssa, \emph{On the Soul and the Resurrection}, trans.\ William Moore, \emph{Nicene and Post-Nicene Fathers}, 2nd series, vol.~5 (New York: Christian Literature Company, 1893), p.~452.}
\end{quote}

Bernard, at the end of his treatise on the love of God, answers it in the language of the Song of Songs, where the Bridegroom bids his friends eat and drink and be inebriated. The fourth degree of love, the love of God for his own sake alone, is reached in its fullness only in the kingdom, and Bernard describes the blessed who have reached it: ``From thence comes satiety which is never satisfied; ardour which is insatiable, yet most calm and peaceful; the eternal and incomparable desire to possess, which arises from no want.''\endnote{Bernard of Clairvaux, \emph{On the Love of God} XI, trans.\ Marianne Caroline and Coventry Patmore, \emph{Saint Bernard on the Love of God}, 2nd ed.\ (London: Burns and Oates, 1884).} Alphonsus, writing for ordinary Christians, says it plainly: the soul ``shall always desire, and always possess that joy. She shall be always satisfied, and always thirsty: for ever thirsty, and always satiated with delights: because the desire of Heaven begets no pain, and its possession produces no tediousness.''\endnote{Alphonsus Liguori, \emph{Preparation for Death}, Consideration XXIX, third point (Duffy, 1843).} The hunger that goes on while it is fed, which both saints describe, is already in Scripture, where Wisdom says of herself in Ecclesiasticus: ``They that eat me, shall yet hunger: and they that drink me, shall yet thirst.''\endnote{Ecclus 24:29, Douay--Rheims (Challoner).}

The Church has left open how this is to be understood. Gregory speaks of an enlargement that never stops; Thomas teaches that each of the blessed reaches the term appointed for him, so that ``nothing further will remain to which he may tend,'' and that heaven is no longer a time ``for advancing by means of merit to a yet greater reward.''\endnote{Aquinas, \emph{Summa theologiae} II-II, q.~28, a.~3, reply to obj.~2; q.~26, a.~13, trans.\ English Dominican Province.} On either view the blessed both rest and are never wearied: they rest because they possess their end, and they are never wearied because that end is the infinite God, whom they never comprehend.

\subsection{The eighth day}

Augustine closed the twenty-two books of \emph{The City of God} with the same rest the \emph{Confessions} had longed for on its first page. God himself, he says, will be the reward of virtue, ``for, as there is nothing greater or better, He has promised Himself.'' And then:

\begin{quote}
He shall be the end of our desires who shall be seen without end, loved without cloy, praised without weariness.\endnote{Augustine, \emph{The City of God} XXII.30, trans.\ Dods (1871).}
\end{quote}

He counts the ages of the world as the days of creation, and sees the saints' rest as a Sabbath that has no evening, ``brought to a close, not by an evening, but by the Lord's day, as an eighth and eternal day, consecrated by the resurrection of Christ, and prefiguring the eternal repose not only of the spirit, but also of the body.'' His last word on the life of heaven is four verbs:

\begin{quote}
There we shall rest and see, see and love, love and praise. This is what shall be in the end without end. For what other end do we propose to ourselves than to attain to the kingdom of which there is no end?\endnote{Augustine, \emph{The City of God} XXII.30, trans.\ Dods (1871).}
\end{quote}

Augustine repeats each verb to begin the next pair, and each step follows from the one before. The restless heart rests because it has found God, and what it has found it sees. To see the supreme Good as he is cannot be anything but to love him, and love that has what it loves overflows in praise. In the same chapter Augustine quotes the psalm: ``Blessed are they that dwell in thy house, O Lord: they shall praise thee for ever and ever.''\endnote{Ps 83:5, Douay--Rheims (Challoner); quoted by Augustine at \emph{The City of God} XXII.30.}

The restlessness we carry now, then, is to be kept and turned toward its end rather than quieted with lesser things, as the psalmist kept it: ``How lovely are thy tabernacles, O Lord of hosts! my soul longeth and fainteth for the courts of the Lord.''\endnote{Ps 83:2--3, Douay--Rheims (Challoner).}
